% GENERATED FILE. Edit canonical lessons, then run npm run generate:books.
% canonical-source-hash: fnv1a64:0c96ff4b5f81f8ff
% canonical-lessons: LA-C437-rudens, LA-C437-funis, LA-C437-armamenta, LA-C437-gubernaculum, LA-C437-remigo

\chapter{Ship's rope, Rope, Tackle, Rudder, To row}
\label{ch:la-ship-s-rope-rope-tackle-rudder-to-row}
\hlchaptermodality{437}

\begin{chapteropening}
\textbf{By the end of this chapter:} \emph{I can talk about ropes, steering and rowing: a ship's rope, a rope, the tackle, the rudder, and rowing.}

Complete the last lesson of chapter 437: I can talk about ropes, steering and rowing: a ship's rope, a rope, the tackle, the rudder, and rowing.
\end{chapteropening}

\section[rudēns, rudentis]{rudēns, rudentis — a ship's rope, a line}

\label{lesson:LA-C437-rudens}



\begin{quote}
Before the new one: say the Latin for a rowing bench, then the Latin for for cargo.
\end{quote}

\subsection*{You'll want to know: rudēns, rudentis}
\textbf{rudēns, rudentis} — \textquotedblleft{}a ship's rope, a line\textquotedblright{}.

\begin{quote}
Nautae rudentēs trahunt et vēlum tollunt.
\end{quote}

\begin{quote}
\emph{The sailors pull the ropes and raise the sail.}
\end{quote}

Third declension, masculine: \emph{rudēns, rudentis m.} The ropes of the rigging.

\begin{cousinweb}
The origin of \emph{rudēns} is not settled: Lewis \& Short call it unknown (ancient writers linked it with \emph{rudō} 'to roar', from the creaking of the ropes), and EtymDB lists it without a parent.
\end{cousinweb}

\begin{grammarlens}[title={What you've built}]
One more word for telling, explaining and describing at B1.
\end{grammarlens}

\subsection*{Your turn}
\begin{itemize}
\raggedright
  \item \emph{Say it:} \emph{rudēns, rudentis}
  \item \emph{Say it:} \emph{rudēns, rudentis}, once more
  \item \emph{Say it:} a sentence of your own with \emph{rudēns, rudentis}
  \item \emph{Recall:} say the Latin for a rowing bench, then the Latin for for cargo, then say \emph{rudēns, rudentis} again
\end{itemize}

\begin{tcolorbox}[breakable,colback=teal!4,colframe=teal!35!black,title={Before you move on}]
What does rudēns, rudentis mean? (\textquotedblleft{}A ship's rope, a line\textquotedblright{}.) Say it once more.
\end{tcolorbox}
\section[fūnis, fūnis]{fūnis, fūnis — a rope, a cord}

\label{lesson:LA-C437-funis}



\begin{quote}
Before the new one: say the Latin for for cargo, then the Latin for a ship's rope.
\end{quote}

\subsection*{You'll want to know: fūnis, fūnis}
\textbf{fūnis, fūnis} — \textquotedblleft{}a rope, a cord\textquotedblright{}.

\begin{quote}
Nāvis fūne ad lītus ligātur.
\end{quote}

\begin{quote}
\emph{The ship is tied to the shore with a rope.}
\end{quote}

Third declension, usually masculine: \emph{fūnis, fūnis m.} English \emph{funicular}, a railway pulled by a rope, comes from it.

\begin{cousinweb}
\emph{fūnis} is inherited and is usually traced to PIE \emph{b\textsuperscript{h}end\textsuperscript{h}-} 'to bind, to tie' (EtymDB); Lewis \& Short also suggest, with a perhaps, the Sanskrit root for binding. English \emph{funicular} comes from its diminutive \emph{fūniculus}.
\end{cousinweb}

\begin{grammarlens}[title={What you've built}]
One more word for telling, explaining and describing at B1.
\end{grammarlens}

\subsection*{Your turn}
\begin{itemize}
\raggedright
  \item \emph{Say it:} \emph{fūnis, fūnis}
  \item \emph{Say it:} \emph{fūnis, fūnis}, once more
  \item \emph{Say it:} a sentence of your own with \emph{fūnis, fūnis}
  \item \emph{Recall:} say the Latin for for cargo, then the Latin for a ship's rope, then say \emph{fūnis, fūnis} again
\end{itemize}

\begin{tcolorbox}[breakable,colback=teal!4,colframe=teal!35!black,title={Before you move on}]
What does fūnis, fūnis mean? (\textquotedblleft{}A rope, a cord\textquotedblright{}.) Say it once more.
\end{tcolorbox}
\section[armāmenta, armāmentōrum]{armāmenta, armāmentōrum — a ship's tackle, the rigging; equipment}

\label{lesson:LA-C437-armamenta}



\begin{quote}
Before the new one: say the Latin for a ship's rope, then the Latin for a rope.
\end{quote}

\subsection*{You'll want to know: armāmenta, armāmentōrum}
\textbf{armāmenta, armāmentōrum} — \textquotedblleft{}a ship's tackle, the rigging; equipment\textquotedblright{}.

\begin{quote}
Procella armāmenta nāvis frangit.
\end{quote}

\begin{quote}
\emph{The storm breaks the ship's rigging.}
\end{quote}

Second declension, neuter plural: \emph{armāmenta, armāmentōrum n.} Masts, sails and ropes together. English \emph{armament} comes from it.

\begin{cousinweb}
\emph{armāmenta} is built on \emph{arma} 'gear, equipment' (Lewis \& Short) with the suffix \emph{-mentum}. English \emph{armament} comes from it.
\end{cousinweb}

\begin{grammarlens}[title={What you've built}]
One more word for telling, explaining and describing at B1.
\end{grammarlens}

\subsection*{Your turn}
\begin{itemize}
\raggedright
  \item \emph{Say it:} \emph{armāmenta, armāmentōrum}
  \item \emph{Say it:} \emph{armāmenta, armāmentōrum}, once more
  \item \emph{Say it:} a sentence of your own with \emph{armāmenta, armāmentōrum}
  \item \emph{Recall:} say the Latin for a ship's rope, then the Latin for a rope, then say \emph{armāmenta, armāmentōrum} again
\end{itemize}

\begin{tcolorbox}[breakable,colback=teal!4,colframe=teal!35!black,title={Before you move on}]
What does armāmenta, armāmentōrum mean? (\textquotedblleft{}A ship's tackle, the rigging; equipment\textquotedblright{}.) Say it once more.
\end{tcolorbox}
\section[gubernāculum, gubernāculī]{gubernāculum, gubernāculī — a rudder, a steering oar}

\label{lesson:LA-C437-gubernaculum}



\begin{quote}
Before the new one: say the Latin for a rope, then the Latin for a ship's tackle.
\end{quote}

\subsection*{You'll want to know: gubernāculum, gubernāculī}
\textbf{gubernāculum, gubernāculī} — \textquotedblleft{}a rudder, a steering oar\textquotedblright{}.

\begin{quote}
Gubernātor gubernāculum tenet et nāvem ad portum dūcit.
\end{quote}

\begin{quote}
\emph{The helmsman holds the rudder and guides the ship to harbour.}
\end{quote}

Second declension, neuter: \emph{gubernāculum, gubernāculī n.} Also the helm of the state. English \emph{govern} comes from the verb \emph{gubernāre}.

\begin{cousinweb}
\emph{gubernāculum} is built on \emph{gubernō} 'to steer' (Lewis \& Short), which was borrowed from Greek \emph{kybernaō}. English \emph{govern} comes from \emph{gubernō} through French.
\end{cousinweb}

\begin{grammarlens}[title={What you've built}]
One more word for telling, explaining and describing at B1.
\end{grammarlens}

\subsection*{Your turn}
\begin{itemize}
\raggedright
  \item \emph{Say it:} \emph{gubernāculum, gubernāculī}
  \item \emph{Say it:} \emph{gubernāculum, gubernāculī}, once more
  \item \emph{Say it:} a sentence of your own with \emph{gubernāculum, gubernāculī}
  \item \emph{Recall:} say the Latin for a rope, then the Latin for a ship's tackle, then say \emph{gubernāculum, gubernāculī} again
\end{itemize}

\begin{tcolorbox}[breakable,colback=teal!4,colframe=teal!35!black,title={Before you move on}]
What does gubernāculum, gubernāculī mean? (\textquotedblleft{}A rudder, a steering oar\textquotedblright{}.) Say it once more.
\end{tcolorbox}
\section[rēmigō, rēmigāre]{rēmigō, rēmigāre — to row}

\label{lesson:LA-C437-remigo}



\begin{quote}
Before the new one: say the Latin for a ship's tackle, then the Latin for a rudder.
\end{quote}

\subsection*{You'll want to know: rēmigō, rēmigāre}
\textbf{rēmigō, rēmigāre} — \textquotedblleft{}to row\textquotedblright{}.

\begin{quote}
Ventus nūllus est; itaque nautae rēmigant.
\end{quote}

\begin{quote}
\emph{There is no wind; so the sailors row.}
\end{quote}

A first-conjugation verb: \emph{rēmigō, rēmigāre}; the perfect is rare. The rower is \emph{rēmex, rēmigis}.

\begin{cousinweb}
\emph{rēmigō} is built from \emph{rēmus} 'oar' and \emph{agō} 'to drive' (EtymDB; Lewis \& Short): to drive the oars.
\end{cousinweb}

\begin{grammarlens}[title={What you've built}]
One more word for telling, explaining and describing at B1.
\end{grammarlens}

\subsection*{Your turn}
\begin{itemize}
\raggedright
  \item \emph{Say it:} \emph{rēmigō, rēmigāre}
  \item \emph{Say it:} \emph{rēmigō, rēmigāre}, once more
  \item \emph{Say it:} a sentence of your own with \emph{rēmigō, rēmigāre}
  \item \emph{Recall:} say the Latin for a ship's tackle, then the Latin for a rudder, then say \emph{rēmigō, rēmigāre} again
\end{itemize}

\begin{tcolorbox}[breakable,colback=teal!4,colframe=teal!35!black,title={Before you move on}]
What does rēmigō, rēmigāre mean? (\textquotedblleft{}To row\textquotedblright{}.) Say it once more.
\end{tcolorbox}
